\subsection{Aufgabe 1 -- Wasserwert}
Um den Wasserwert zu bestimmen, benötigt man \(c_{\text{W}}(\qty{50}{\degreeCelsius}) = \qty{4,186(4)}{\J\per\g\per\K}\)\autocite{Skript_1}. Aus \cref{fig:Aufgabe1.pdf} wurde extrapoliert:
\[\tcbhighmath[colback=white]{\overline{T} = \qty{47,0(4)}{\degreeCelsius}}\]
Der Fehler ergibt sich aus der nicht exakt zum Zeitpunkt des Befüllens gestarteten Messreihe. Zusätzlich wurde mit einem ungenauen Wandthermometer die Raumtemperatur \(T_2\) und die Temperatur des Wassers \(T_1\) mit einem Digitalthermometer gemessen:
\begin{align*}
  T_2 = \qty{25,1(5)}{\degreeCelsius} && T_1 = \qty{47,90(14)}{\degreeCelsius}
\end{align*}
dessen systematischer Fehler für diese Aufgabe vernachlässigt wurde, da \(\overline{T}\) und \(T_1\) beide mit dem Digitalthermometer gemessen wurden, und sich somit der systematische Fehler in der Differenz des Zählers von \cref{eq:Wasserwert} (siehe auch unten herauskürzt). Die Masse des Wassers bestimmt sich durch Wiegen des Kalorimeters mit und ohne Wasser sowie eine Differenzbildung zu:
\begin{align}
  m_W=m_1-m_2&&\Delta m = \sqrt{\Delta m_1^{2} + (-\Delta m_2^2)}&&m_W = \qty{224,80(14)}{\g}
\end{align}
Damit ergibt sich mit der Fehlerfortpflanzung von:
\begin{gather}
      W=c_W\,m_W\dfrac{T_1-\overline{T}}{\overline{T}-T_2}\\
      \begin{split}
        \Delta W^2 = \left(\frac{m_{\text{W}} \left(T_1-\overline{T}\right)}{-T_2+\overline{T}} \Delta c_{\text{W}}\right)^2 + \left(\frac{c_{\text{W}} \left(T_1-\overline{T}\right)}{-T_2+\overline{T}} \Delta m_{\text{W}}\right)^2
        + \left(\frac{c_{\text{W}} m_{\text{W}}}{-T_2+\overline{T}} \Delta T_1\right)^2\\ + \left(\frac{c_{\text{W}} m_{\text{W}} \left(T_1-\overline{T}\right)}{\left(-T_2+\overline{T}\right)^2} \Delta T_2\right)^2
        + \left(\left(\frac{-c_{\text{W}} m_{\text{W}} \left(T_1-\overline{T}\right)}{\left(-T_2+\overline{T}\right)^2} - \frac{c_{\text{W}} m_{\text{W}}}{-T_2+\overline{T}}\right) \Delta\overline{T}\right)^2
      \end{split}
\end{gather}
\[\tcbhighmath{W_{\text{exp}} = \qty{39(19)}{\J\per\K}}\]
Dieses Ergebnis weicht absolut sehr stark vom Herstellerwert \(W_{\text{lit}} = \qty{70}{\J\per\K}\) ab: 
\begin{table}[htbp]
\centering
\caption{Vergleich von \(W_{\text{exp}}\) und \(W_{\text{lit}}\)}
\begin{tabularx}{\textwidth}{LZZZx}
    \toprule
        W_{\text{exp}}\;[\unit{\J\per\K}]&W_{\text{lit}}\;[\unit{\J\per\K}]&\text{abs. Abw.}\;[\unit{\J\per\K}]&\text{proz. Abw.}\;[\unit{\percent}]&S\;[\sigma]\\\midrule
        \num{39(19)}&70&\num{31(19)}&\num{80(70)}&\num{1.7}\\
    \bottomrule
\end{tabularx}
\label{table:Vergleich_W} 
\end{table}

\xfigpdf{p}{width=\textwidth}{Aufgabe1.pdf}{}{Grafische Ermittlung von \(\overline{T}\)}
\clearpage

\subsection[\texorpdfstring{Aufgabe 2 -- Bestimmung von \(c_{\mathrm{mol}}\) der Probekörper}{Aufgabe 2 -- Bestimmung von cmol der Probekörper}]{Aufgabe 2 -- Bestimmung von \(\boldsymbol{c_{\mathrm{mol}}}\) der Probekörper}
Mit \cref{c_x} sowie der Einbeziehung, dass Körper 2 aus dieser Gleichung aus einem wassergefüllten Kalorimeter besteht (also aus zwei Bestandteilen), ergibt sich für die spezifische Wärmekapazität \(c_x\) der Probekörper:
\begin{equation}
    c_x = \frac{(m_{\text{W}} c_{\text{W}} + W)(\overline{T} - T_2)}{m_x (T_1 - \overline{T})}
    \label{c_x_kalorimeter}
\end{equation}
\begin{equation}
    \begin{aligned}
        \Delta c_x^{2} = &\left(\frac{-\left(\overline{T}-T_2\right) \left(W+c_{\text{W}} m_{\text{W}}\right)}{{m_x}^2 \left(-\overline{T}+T_1\right)} \Delta m_x\right)^2 + \left(\frac{c_{\text{W}}(\overline{T}-T_2)}{m_x (T_1-\overline{T})}\Delta m_{\text{W}}\right)^2 \\
        + &\left(\frac{m_{\text{W}}(\overline{T}-T_2)}{m_x (T_1-\overline{T})}\Delta c_{\text{W}}\right)^2 + \left(\frac{\overline{T}-T_2}{m_x \left(-\overline{T}+T_1\right)} \Delta W\right)^2 \\
        + &\left(\frac{-\left(\overline{T}-T_2\right) \left(W+c_{\text{W}} m_{\text{W}}\right)}{m_x \left(-\overline{T}+T_1\right)^2} \Delta T_1\right)^2 + \left(\frac{-\left(W+c_{\text{W}} m_{\text{W}}\right)}{m_x \left(-\overline{T}+T_1\right)} \Delta T_2\right)^2 \\
        + &\left(\left(\frac{W+c_{\text{W}} m_{\text{W}}}{m_x \left(-\overline{T}+T_1\right)} + \frac{\left(\overline{T}-T_2\right) \left(W+c_{\text{W}} m_{\text{W}}\right)}{m_x \left(-\overline{T}+T_1\right)^2}\right) \Delta \overline{T}\right)^2
    \end{aligned}
    \label{c_x_fehler}
\end{equation}
Die Temperatur \(T_1\) der Probekörper beträgt ungefähr \qty{100}{\degreeCelsius}, da sie in siedendes Wasser getaucht wurden. Die Siedetemperatur ist jedoch luftdruckabhängig; gemessen bei \(p = \qty{1001,8(10)}{\hPa}\) mit:
\begin{equation}
    T_1 = \qty{100}{\degreeCelsius} + \qty{0,0276}{\degreeCelsius\per\hPa} \cdot \left(p - \qty{1013}{\hPa}\right)
    \label{luftdruck}
\end{equation}
\begin{equation}
    \Delta T_1 = \num{0,0276} \cdot \Delta p
\end{equation}
ergibt sich:
\[T_1 = \qty{99,690(28)}{\degreeCelsius}\]
Aus den Tabellen im Messprotokoll wurden jeweils die Massen \(m_x\) der 3 Probekörper, die Masse des Wassers \(m_{\text{W}}\) im Kalorimeter (welche durch das Ein- und Ausbringen sinkt) mit gleichbleibendem Fehler \(\Delta m_{\text{W}} = \qty{0,14}{\g}\), die Mischtemperatur (immer der maximale Wert) sowie die Temperatur des Kalorimeters \(T_2\) abgelesen. Der systematische Fehler von \(T_2\) könnte weiterhin vernachlässigt werden, da \(\overline{T}\) und \(T_2\) beide mit dem Digitalthermometer gemessen wurden, \(T_1\) hingegen nicht, wodurch sich der systematische Fehler in der Nenner-Differenz \(T_1-\overline{T}\) nicht herauskürzen würde. Da die Mischtemperatur weiterhin \(<\qty{65}{\degreeCelsius}\) ist (dem Bereich für die Gültigkeit des gegebenen \(c_W\) Wertes). bleibt \(c_W\) bleibt beim selben Wert; der Wasserwert wurde nun dem Skript mit pauschal gesetztem Fehler entnommen: \(W = \qty{70(10)}{\J\per\K}\).

\renewcommand{\arraystretch}{1.3}  
\begin{table}[htbp]
\centering
\caption{Werte aus dem Messprotokoll mit Fehlern}
\begin{tabularx}{\textwidth}{lZZZZZZ}
\toprule
Probekörper & m_x \;[\unit{\g}] & m_{\text{W}} \;[\unit{\g}] & \overline{T} \;[\unit{\degreeCelsius}] & T_2 \;\;[\unit{\degreeCelsius}] & W \;[\unit{\J\per\K}] \\ 
\midrule
Graphit & \num{125,2(14)} & \num{348,9(14)} & \num{29,1(11)} & \num{24,40(07)} & \num{70(10)} \\
Blei & \num{682,7(14)} & \num{342,6(14)} & \num{32,5(11)} & \num{28,80(09)} & \num{70(10)} \\
Aluminium & \num{173,7(14)} & \num{341,5(14)} & \num{38,0(11)} & \num{32,20(10)} & \num{70(10)} \\
\bottomrule
\end{tabularx}
\label{table:Werte_c_x}
\end{table}
\FloatBarrier

Als Ergebnisse für \(c_x\) und \(c_{\text{mol}}=cM\) berechnen sich mithilfe der Literaturwerte der molaren Masse \(M_x\) aus \cref{table:Literatur} und \(\Delta c_{\text{mol}} = M \cdot \Delta c\) die Werte in der folgenden Tabelle. Die Umrechnung der Temperaturen in Kelvin kann man sich sparen, da in den Gleichungen immer nur die Differenz zweier Temperaturen auftaucht; die Kelvin-Skala ist nur verschoben, das kürzt sich heraus.
\begin{table}[htbp]
\centering
\caption{Spezifische Wärmekapazität und Molwärme bei der Wassermessung}
\begin{tabularx}{0.8\textwidth}{TZZ}
\toprule
Material &c_x \;[\unit{\J\per\g\per\K}] & c_{x,\text{mol}} \;[\unit{\J\per\mol\per\K}] \\
\midrule
Graphit & \num{0,81(27)} & \num{10(4)} \\
Blei & \num{0,12(5)} & \num{25(13)} \\
Aluminium & \num{0,81(23)} & \num{22(7)} \\
\bottomrule
\end{tabularx}
\label{table:c_x_wasser}
\end{table}

\paragraph{Vergleich mit Literaturwerten}
Es werden nun die gemessenen \(c_x\) mit den Literaturwerten sowie die gemessenen Werte mit den theoretisch nach \cref{Dulong_x} bestimmten \(c_{x,\text{Dulong}}\) verglichen.

\begin{table}[htbp]
\centering
\caption{Abweichungen der \(c_x\) der Wassermessung zu Literatur- und Dulong-Werten}
\begin{tabularx}{\textwidth}{lZZS[table-format=1.2]ZS[table-format=1.2]}
\toprule
Material & c_{x,\text{Lit.}} \;[\unit{\J\per\g\per\K}] & c_{x,\text{gem.}} \;[\unit{\J\per\g\per\K}] & {\(S\;[\sigma]\)} & c_{x,\text{Dulong}} \;[\unit{\J\per\g\per\K}] & {\(S\;[\sigma]\)} \\
\midrule
Graphit & \num{0,709} & \num{0,81(27)} & \num{0,37} & \num{2,077} & \num{4,7} \\
Blei & \num{0,129} & \num{0,13(6)} & \num{0,17} & \num{0,120} & \num{0,16} \\
Aluminium & \num{0,90} & \num{0,81(23)} & \num{0,4} & \num{0,924} & \num{0,5} \\
\bottomrule
\end{tabularx}
\label{table:c_x_vergleich}
\end{table}

\newpage
\subsection[\texorpdfstring{Aufgabe 3 -- \(c_{\mathrm{mol}}\) über Stickstoffverdampfung}{Aufgabe 3 -- cmol über Stickstoffverdampfung}]{Aufgabe 3 -- \(\boldsymbol{c_{\mathrm{mol}}}\) über Stickstoffverdampfung}
Mit \cref{c_x_Stickstoff} und der Fehlerfortpflanzung:
\begin{equation}
    \Delta c_{x,\ce{N2}}^2 = \left(\frac{-Q_V m_V}{{m_x}^2 \left(-T_{\ce{N2}}+T_1\right)} \Delta m_x\right)^2 + \left(\frac{Q_V}{m_x \left(-T_{\ce{N2}}+T_1\right)} \Delta m_V\right)^2 + \left(\frac{-Q_V m_V}{m_x \left(-T_{\ce{N2}}+T_1\right)^2} \Delta T_1\right)^2
\end{equation}
sowie dem Wert \(Q_V = \qty{199}{\J\per\g}\), den Werten aus dem Messprotokoll und \(T_{\ce{N2}} = \qty{-195,8}{\degreeCelsius}\) ergeben sich:

\begin{table}[htbp]
\centering
\caption{Spezifische Wärmekapazität und Molwärme bei der Stickstoffmessung}
\begin{tabularx}{0.8\textwidth}{TZZ}
\toprule
Material & c_{x,\ce{N2}} \;[\unit{\J\per\g\per\K}] & c_{x,\text{mol}}(\ce{N2}) \;[\unit{\J\per\mol\per\K}] \\
\midrule
Graphit & \num{0,467(4)} & \num{5,61(5)} \\
Blei & \num{0,1322(11)} & \num{27,39(23)} \\
Aluminium & \num{0,719(6)} & \num{19,40(16)} \\
\bottomrule
\end{tabularx}
\label{table:c_x_N}
\end{table}

Wir haben hier, wie eingangs erwähnt, nur den Mittelwert der spezifischen Wärmekapazität bzw. Molwärme aus den zwei Temperaturbereichen bestimmt. In \cref{fig:c_T.png} sieht man, dass diese Mittelwertbildung für den roten Bereich (Messung mit siedendem Wasser) eine gute Näherung bietet, da \(c_{x,\text{mol}}\) nicht sehr stark schwankt. Für den schwarzen Bereich (Messung mit flüssigem Stickstoff) ist die Mittelwertbildung keine gute Näherung mehr, da sich \(c_{x,\text{mol}}\) in diesem Bereich stark ändert.

Aus diesem Grund sind auch nur die Literaturwerte für \(c_x(\qty{300}{\K})\) angegeben. Natürlich hätte im Skript auch der Mittelwert \(\overline{c_x}(\qty{77}{\K} \dots \qty{300}{\K})\) angegeben werden können; unsere Raumtemperatur beträgt jedoch auch nicht exakt \qty{300}{\K}, was allerdings vernachlässigbar wäre.

\newpage
\subsection{Aufgabe 4 -- Verhältnisse}
Das in \cref{fig:bild2} dargestellte Verhältnis \(V \defeq \frac{c_{\text{mol}}(\ce{N2})}{c_{\text{mol}}(\ce{H2O})}\) soll berechnet werden, um durch Ablesen in derselben Abbildung die Debye-Temperatur der Stoffe zu bestimmen. Der Fehler des Verhältnisses bestimmt sich zu:
\begin{equation}
    \Delta V = \sqrt{\left(\frac{1}{c_{\ce{H2O}}} \Delta c_{\ce{N2}}\right)^2 + \left(\frac{-c_{\ce{N2}}}{c_{\ce{H2O}}^2} \Delta c_{\ce{H2O}}\right)^2}
\end{equation}

\begin{table}[htbp]
\centering
\caption{Vergleich der Molwärmen der beiden Messungen und Bestimmung von \(\Theta_D\)}
\begin{tabularx}{\textwidth}{Txxxxx}
\toprule
\rule{0pt}{2.5ex}Material & c_{x,\text{mol}}(\ce{H2O}) \;[\unit{\J\per\mol\per\K}] & c_{x,\text{mol}}(\ce{N2}) \;[\unit{\J\per\mol\per\K}] & V & \Theta_{D,\text{abgel.}} \;[\unit{\K}] & \Theta_{D,\text{Lit.}} [\unit{\K}] \\
\midrule
Graphit & \num{10(4)} & \num{5,61(5)} & \num{0,58(19)} & \num{800(500)} & 2230 \\
Blei & \num{25(13)} & \num{27,39(23)} & \num{1,1(5)} & \num{0(1000)} & 95 \\
Aluminium & \num{22(7)} & \num{19,40(16)} & \num{0,89(25)} & \num{300(350)} & 430 \\
\bottomrule
\end{tabularx}
\label{table:theta_D}
\end{table}